% Omar - Ejercicio 1: Preguntas de repaso.

\section{Preguntas de repaso}

Las respuestas siguientes se apoyan en una distinción que atraviesa todo el
modelo relacional: el \emph{esquema} describe la estructura de los datos y el
\emph{estado} describe el conjunto de hechos que, en un momento dado,
satisface esa estructura. Confundir ambos niveles lleva a preguntarse, por
ejemplo, si una llave existe en una base vacía, pregunta que sólo tiene
sentido en el primer nivel \cite{elmasri2016fundamentals}.

\subsection{La relación y sus características}

Una \textbf{relación} es un conjunto de tuplas homogéneas que representa
hechos acerca de una entidad o de una asociación entre entidades. En términos
precisos, dado un nombre de relación $R$ y un conjunto de dominios $D_1, D_2,
\ldots, D_n$, la relación es un subconjunto del producto cartesiano

\begin{equation}
  R \subseteq D_1 \times D_2 \times \cdots \times D_n,
\end{equation}

de modo que cada elemento de $R$ es una $n$-tupla $(a_1, a_2, \ldots, a_n)$
con $a_i \in D_i$. La interpretación usual es que cada posición de la tupla
corresponde a un atributo con nombre propio; por ello dos tuplas que coinciden
componente a componente se consideran la misma tupla, aunque se hayan escrito
en distinto orden.

La noción de relación es más restrictiva que la de tabla, y sus
características son precisamente las que la hacen útil como base de un modelo
formal \cite{codd1970relational}. Las seis propiedades siguientes enuncian el
carácter de conjunto y de estructura fija que debe tener una relación:

\begin{enumerate}[leftmargin=1.9em,itemsep=0.4em]
  \item \textbf{Es un conjunto.} No puede contener tuplas repetidas. Si dos
    filas tienen exactamente los mismos valores, son la misma tupla y la
    segunda no se almacena. Es la característica que más claramente distingue
    una relación de una lista, donde la repetición y el orden sí tienen
    significado.
  \item \textbf{Las tuplas son homogéneas.} Todas tienen el mismo número de
    componentes y cada componente ocupa una posición con un dominio asociado.
    No existen tuplas de menor o mayor longitud dentro de una misma relación.
  \item \textbf{Los componentes están ordenados y son atómicos.} El orden de
    los atributos es parte de la estructura del esquema y se establece al
    declarar la relación; además, cada componente corresponde a un único valor
    del dominio y no a un conjunto de valores. La consecuencia de la
    atomicidad es que un atributo compuesto no puede ocupar una sola posición:
    sus partes deben ocupar posiciones distintas.
  \item \textbf{El orden de las tuplas es irrelevante.} La relación es un
    conjunto y, en consecuencia, no admite un orden canónico. Recuperar las
    filas en un orden distinto no produce una relación distinta.
  \item \textbf{Cada componente pertenece a un dominio.} El dominio acota los
    valores admisibles y aporta el tipo y el rango, lo que impide mezclar
    incoherentemente, por ejemplo, un número de empleado con un nombre.
  \item \textbf{La instancia y el tipo son objetos distintos.} El tipo de
    relación, formado por el nombre, el grado y los dominios, pertenece al
    esquema; la relación como subconjunto del producto cartesiano es una
    instancia, es decir, parte del estado y susceptible de cambiar.
\end{enumerate}

\begin{conceptobox}[Consecuencia práctica]
Que una relación sea un conjunto implica que la búsqueda de una tupla es una
  búsqueda de miembro y no de posición. Ésta es la razón formal de que un
  modelo relacional no pueda ofrecer la noción de última fila insertada: el
  orden en que las instancias se incorporan al estado no forma parte de la
  relación.
\end{conceptobox}

\subsection{El esquema de relación}

El \textbf{esquema de relación} es la definición del tipo de relación: es la
especificación de la estructura contra la que se valida cualquier instancia.
Se enuncia como un nombre seguido de los atributos con su dominio

\begin{equation}
  R(A_1{:}D_1,\; A_2{:}D_2,\; \ldots,\; A_n{:}D_n),
\end{equation}

junto con las restricciones que el modelo impone sobre la relación. La parte
relacionada con la estructura ---el nombre, el número de atributos y sus
dominios--- se denomina \emph{intensión} de la relación, en tanto que el valor
concreto que la relación adopta en un estado determinado es su
\emph{extensión}.

Cabe subrayar que un esquema relacional no queda determinado por la lista de
atributos. Cuando un modelo relacional se obtiene por traducción de un modelo
Entidad--Relación, el esquema incorpora además las llaves primarias, las
llaves foráneas, las dependencias funcionales y las restricciones de
integridad que la traducción haya hecho explícitas. En el sublenguaje SQL esta
correspondencia es directa: el nombre de la tabla y sus columnas son la
intensión, y las cláusulas \texttt{PRIMARY KEY}, \texttt{FOREIGN KEY},
\texttt{UNIQUE} y \texttt{CHECK} son la parte declarativa de las
restricciones. Un esquema que omite tales restricciones admite instancias que
satisfacen la declaración pero que son semánticamente imposibles en el
mini--mundo que se quiere representar.

\subsection{Superllave, llave mínima, llave candidata y llave primaria}

Las cuatro nociones se ordenan por inclusión estricta: toda llave primaria es
una llave candidata, toda llave candidata es una llave mínima y toda llave
mínima es una superllave. La diferencia entre ellas no está en lo que
identifican, sino en el grado de justificación que exigen
\cite{elmasri2016fundamentals}.

\medskip
\noindent\textbf{Superllave.} Conjunto de atributos $K$ tal que dos tuplas
  distintasde la relación no pueden coincidir en todos los atributos de $K$.
  Equivalentemente,$K$ determina funcionalmente al resto de los atributos de
  larelación. La condición es suficiente pero no necesaria: se admiten
  superllavesque contienen atributos innecesarios, por lo que la definición
  noimpone economía alguna. Todo conjunto que contenga una llave es,
  automáticamente,una superllave.
\medskip
\noindent\textbf{Llave mínima.} Superllave de la que no puede extraerse
  ningúnatributo sin perder la propiedad de superllave. Equivale a exigir que
  todosubconjunto propio de $K$ deje de ser superllave. La minimalidad es una
  propiedadque se comprueba atributo por atributo y sólo tiene sentido
  respectode una superllave determinada.
\medskip
\noindent\textbf{Llave candidata.} Llave mínima, esto es, una superllave sin
  atributosredundantes. Es el candidato más económico que el modelo ofrece
  paraidentificar las tuplas de la relación. Una relación puede tener varias
  llavescandidatas, y el esquema debe documentar todas ellas, incluso si sólo
  unase designa como primaria: las restantes no dejan de ser candidatas por
  nohaber sido elegidas.
\medskip
\noindent\textbf{Llave primaria.} La llave candidata que el diseñador elige
  comoidentificador único de las tuplas de la relación. Se declara
  explícitamenteen el esquema, sus atributos pasan a ser obligatorios y no
  admitenulos, y su valor es el que las demás relaciones referencian mediante
  llavesforáneas. Existe a lo sumo una llave primaria por relación; elegida
  ella,las demás llaves candidatas se reducen a candidatas alternativas o se
  declinanmediante la restricción \texttt{UNIQUE}.

La diferencia entre candidata y primaria suele originar confusión. Que un
atributo no sea primario no significa que sea irrelevante para la
identificación: si también determina al resto de los atributos, sigue siendo
candidato. Lo que hace la llave primaria es notificar la elección en el
esquema, y esa notificación es la que las demás relaciones aprovechan al
declarar sus llaves foráneas.

Para ilustrar las cuatro nociones con un esquema de empleados, considérese

\begin{equation}
  \mathit{Empleado}(\mathit{N\acute{u}mEmpleado},\;
  \mathit{Nombre},\; \mathit{Paterno},\; \mathit{Materno},\;
  \mathit{Nacimiento},\; \mathit{G\acute{e}nero},\;
  \mathit{FechaContrataci\acute{o}n}).
\end{equation}

De $\{N\acute{u}mEmpleado\}$ se afirma en el esquema que es llave primaria y
por tanto candidata y mínima. El conjunto completo de los siete atributos es
también superllave, pero no es mínima: al retirar \emph{Paterno},
\emph{Materno}, \emph{Nacimiento}, \emph{G\'enero} y
\emph{FechaContrataci\'on} se conserva la propiedad, y lo mismo ocurre al
retirar únicamente $N\acute{u}mEmpleado$. La combinación $\{Nombre, Paterno,
Materno, Nacimiento\}$ es superllave en el dominio del negocio y, en la
mayoría de los casos, también es mínima, porque ninguno de sus cuatro
atributos resulta superfluo; se trata por tanto de una llave candidata que el
esquema no eligió como primaria. La omisión de esa declaración se documenta
como una decisión de diseño: se descarta la dependencia de los nombres y de la
fecha de nacimiento por su propia inestabilidad ante la homonimia y ante los
cambios de nombre.

\subsection{Restricciones que imponen las llaves primaria y foránea}

Una llave primaria impone un conjunto de restricciones que sólo se enuncia
respecto del esquema y que el sistema debe hacer cumplir en todo momento:

\begin{itemize}
  \item \textbf{Unicidad.} Dos tuplas distintas de la relación no pueden
    coincidir en el valor de la llave primaria. Es la definición de
    superllave, aplicada a los atributos elegidos como primarios.
  \item \textbf{No nulabilidad.} Ningún componente de la llave primaria puede
    tomar el valor nulo. Esta restricción es lo que garantiza que la llave
    exista siempre y por tanto identifique siempre a la tupla; sin ella, la
    noción de identificador se vuelve ambigua y las referencias desde otras
    relaciones quedan sin destino.
  \item \textbf{Identificación.} La llave primaria determina funcionalmente a
    todos los demás atributos de la relación, lo que permite afirmar que una
    entidad del mini--mundo queda descrita por los datos de una sola tupla.
  \item \textbf{Existencia de un punto de referencia.} Al ser el único
    mecanismo de direccionamiento entre relaciones, su declaración convierte a
    la relación en el destino válido de las llaves foráneas que otras
    relaciones declaren.
\end{itemize}

Una llave foránea impone, por su parte, la restricción de integridad
referencial, que se apoya en la compatibilidad entre el atributo o conjunto de
atributos declarados y la llave primaria de la relación referenciada:

\begin{itemize}
  \item \textbf{Integridad de la referencia.} El valor de la llave foránea
    debe coincidir con el valor de la llave primaria de alguna tupla de la
    relación referenciada. Si el valor es nulo y la columna lo permite, la
    referencia queda simplemente sin destino; en caso contrario, la inserción
    debe rechazarse.
  \item \textbf{Tipos compatibles.} El número de atributos, su orden y sus
    dominios deben corresponder con los de la llave primaria referenciada. Una
    llave foránea sobre columnas de tipo distinto no puede validarse y debe
    rechazarse en el esquema, no en tiempo de ejecución.
  \item \textbf{Efecto sobre la operación.} La restricción condiciona el orden
    de las operaciones: un hijo no puede insertarse antes que su padre, y un
    padre no puede eliminarse mientras tener hijos que lo referencien, salvo
    que la definición de la llave foránea especifique una acción determinada.
    Las acciones habituales son restringir la operación, propagar en cascada o
    asignar nulo a la referencia.
  \item \textbf{No confunde llave foránea con participación total.} Que el
    lado foráneo sea obligatorio y que el lado referenciado admita varias
    referencias son hechos distintos. La llave foránea dice que cada tupla
    hija apunta exactamente a un padre; no dice que todo padre deba tener
    hijos.
  \item \textbf{Puede formar parte de la llave primaria y, entonces, no ser
    nulable.} Cuando un atributo es simultáneamente parte de la llave primaria
    de su relación y llave foránea de otra, hereda la no nulabilidad de la
    primera. Es el caso típico de las relaciones que registran hechos
    asociados a un periodo, como las que aparecen en el Ejercicio~4.
\end{itemize}

\begin{conceptobox}[Distinción operativa]
La llave primaria protege la \emph{identidad} de la tupla dentro de su
  relación; la llave foránea protege la \emph{existencia} de lo que la tupla
  referencia. La primera hace imposible tener dos tuplas iguales; la segunda
  hace imposible tener una tupla que apunte a la nada.
\end{conceptobox}

\subsection{Las reglas de Codd}

En 1970, Edgar F. Codd publicó un modelo para el manejo de grandes bases de
datos compartidas y estableció un conjunto de criterios que permitían
distinguir un sistema relacional de un sistema de archivos organizado en
jerarquías \cite{codd1970relational}. Esas doce reglas enumeran propiedades
que un sistema debe satisfacer para poder considerarse relacional; no
constituyen una especificación técnica, sino un criterio de diseño que sigue
siendo el marco habitual para discutir la calidad de un sistema de gestión de
bases de datos \cite{elmasri2016fundamentals}.

Las ocho primeras reglas se refieren a una base de datos personal, es decir,
no distribuida, y las cuatro restantes se ocupan de la distribución.

\subsubsection*{Reglas de acceso a los datos}

\begin{enumerate}[leftmargin=1.9em,itemsep=0.55em]
  \item \textbf{Regla de información.} Todo dato almacenado debe poder
    localizarse de manera unívoca, y únicamente, con el nombre de la relación,
    el nombre del atributo y el valor de la llave. La regla prohíbe, en
    consecuencia, direccionar los datos mediante rutas físicas, direcciones de
    disco o punteros. En términos de los apartados anteriores, exige que todo
    hecho esté en una tupla y que esa tupla sea alcanzable sin conocer la
    organización del almacenamiento.
  \item \textbf{Regla de acceso garantizado.} Si un usuario tiene autorizada
    la consulta, el sistema debe permitirle acceder a cualquier dato de la
    base. No puede existir información que sea alcanzable un día e
    inalcanzable al siguiente. La consecuencia práctica es que el esquema debe
    reflejar toda la información que el mini--mundo requiere; si un dato no
    puede nombrarse, es un dato mal modelado.
\end{enumerate}

\subsubsection*{Reglas de integridad de los datos}

\begin{enumerate}[leftmargin=1.9em,itemsep=0.55em,start=3]
  \item \textbf{Regla de tratamiento sistemático de los valores nulos.} La
    información ausente debe representarse mediante un mecanismo explícito y
    homogéneo. El modelo relacional original no admite el valor nulo como
    miembro de un dominio y prefiere, para cada situación, o bien separar el
    atributo opcional en otra relación, o bien dividir la relación en dos. La
    regla admite cualquier método sistemático, pero exige que sea uniforme y
    conocido, porque de lo contrario la semántica de la base queda
    indeterminada.
  \item \textbf{Regla de datos dinámicos.} El contenido de la base puede
    cambiar sin que sea necesario modificar su definición. Insertar o eliminar
    una tupla es una operación ordinaria del modelo y no obliga a redefinir el
    esquema, a diferencia de lo que ocurría en los sistemas de archivos, donde
    añadir un campo obligaba a reorganizar los registros.
  \item \textbf{Regla de datos consolidados.} Un hecho determinado debe
    almacenarse en un solo lugar de la base de datos, y todos los datos
    necesarios para identificarlo y describirlo deben residir en esa misma
    tupla o en tuplas relacionadas de manera explícita. La regla admite la
    redundancia, que es la base de la normalización, pero exige que la
    duplicación no contenga copias desactualizadas: si un dato se repite,
    todas sus copias deben coincidir.
  \item \textbf{Regla de almacenamiento en forma lógica.} El usuario accede a
    los datos por su significado y no por su ubicación física. El conocimiento
    de los nombres de relación y de atributo debe bastar para operar sobre la
    base, sin pedirle que se entienda cómo están organizados los registros en
    disco.
  \item \textbf{Regla de independencia física de los datos.} La organización
    física de los datos puede modificarse sin afectar a las aplicaciones ni a
    las consultas lógicas. Es la regla que autoriza añadir un índice,
    reorganizar los archivos o repartir los registros en distintos
    dispositivos de almacenamiento sin reescribir el código que los consulta.
  \item \textbf{Regla de independencia lógica de los datos.} La estructura
    lógica de los datos puede modificarse de una manera que no afecte a las
    aplicaciones ni a sus datos. A diferencia de la anterior, aquí los cambios
    --añadir una relación, dividir una en dos, reorganizar los atributos-- sí
    modifican la forma en que se presenta la información, y la regla exige que
    ese cambio pueda realizarse sin la intervención de los programas.
\end{enumerate}

\subsubsection*{Reglas de distribución de los datos}

\begin{enumerate}[leftmargin=1.9em,itemsep=0.55em,start=9]
  \item \textbf{Regla de distribución física.} La base de datos puede estar
    repartida físicamente entre varios sitios de almacenamiento, por ejemplo
    en una red de servidores o en dispositivos distintos, y el usuario debe
    poder consumirla como si fuera una sola base centralizada.
  \item \textbf{Regla de distribución lógica.} La base de datos sigue siendo
    relacional para el usuario aunque su implementación se encuentre
    físicamente repartida. La experiencia de uso --consultar y actualizar
    relaciones completas mediante un lenguaje declarativo-- no cambia por el
    hecho de que los datos residan en máquinas distintas.
\end{enumerate}

\subsubsection*{Reglas de sustitutibilidad}

\begin{enumerate}[leftmargin=1.9em,itemsep=0.55em,start=11]
  \item \textbf{Regla de sustitutibilidad de los datos.} Si un usuario opera
    sobre una relación derivada, por ejemplo una vista, debe disponer de las
    mismas oportunidades de lectura y de escritura que si operara sobre la
    relación almacenada. En términos del modelo, una vista es una relación,
    con tuplas y, por extensión, con instancias, inasmuch como cualquier otra.
  \item \textbf{Regla de sustitutibilidad de la integridad.} Las restricciones
    del esquema de la base almacenada, incluidas sus llaves primarias y
    candidatas, deben conservarse en cualquier relación derivada. Por
    consiguiente, los datos derivados tienen la misma fiabilidad que los
    almacenados, y una vista jamás puede revelar un estado que la base real no
    permita.
\end{enumerate}

\begin{conceptobox}[Alcance de las doce reglas]
Las reglas del bloque central remiten directamente a las cuatro llaves
  estudiadas en el apartado anterior. La regla de información supone la
  existencia de llaves primarias que permitan designar toda tupla; la regla de
  integridad exige que no se declaren nulos donde la llave primaria ya los
  prohíbe; la regla de datos consolidados remite a la noción de superllave, ya
  que un hecho se identifica por una combinación mínima de atributos; y la
  regla de sustitutibilidad de la integridad exige que las relaciones
  derivadas respeten las mismas claves que las relaciones almacenadas.
\end{conceptobox}

\subsubsection*{Por qué son importantes}

Las reglas de Codd son importantes porque desplazan el objeto de la
construcción: no un programa que accede a sus datos, sino un sistema cuya
descripción lógica de los datos constituye el punto de partida del trabajo.

\begin{itemize}
  \item \textbf{Hacen posible la independencia de los datos.} Las reglas siete
    y ocho convierten la independencia física y lógica en un requisito del
    modelo y no en un accidento de la implementación. Gracias a ellas, una
    aplicación puede permanecer escrita una sola vez mientras la base cambia
    de organización, y un mismo esquema puede servir a aplicaciones distintas
    sin que cada una conozca la representación interna. Ésta es,
    probablemente, la consecuencia práctica de mayor alcance de todo el
    modelo.
  \item \textbf{Sostienen los lenguajes declarativos de consulta.} Si la regla
    de almacenamiento en forma lógica obliga a acceder a los datos por su
    significado y no por su recorrido físico, entonces el usuario puede
    describir qué desea obtener en lugar de cómo obtenerlo. Ésa es la
    condición precisa para que exista un optimizador que elija el plan de
    ejecución. En la práctica, el álgebra relacional y el SQL son la
    manifestación de esta regla, y el optimizador de consultas es la prueba de
    que el modelo tuvo efecto.
  \item \textbf{Dan fundamento a las vistas.} Las reglas once y doce autorizan
    a tratar una relación derivada como si fuera almacenada, en igualdad de
    condiciones para leer y para escribir. Ésa es la razón por la que la
    lógica de vistas de SQL es exacta en lugar de aproximada: la vista no es
    una copia que pueda desactualizarse, sino una fórmula sobre el estado
    real.
  \item \textbf{Fundamentan el tratamiento de la integridad.} La regla de
    datos consolidados obliga a identificar dónde reside cada hecho y, con
    ello, a determinar qué atributos dependen de cuáles. Ese análisis de
    dependencias funcionales es el fundamento de la normalización y de la
    detección de llaves candidatas, de modo que las reglas conectan la
    formulación relacional con el diseño de esquemas que se estudia al aplicar
    el modelo relacional.
  \item \textbf{Anticipan las bases de datos distribuidas.} Las reglas nueve a
    doce extienden la formulación inicial a un contexto en el que los datos
    residen en varios servidores. La anticipación es notable, porque fue
    enunciada en 1970, mucho antes de que la distribución se generalizara en
    la práctica. El hecho de que las reglas resulten anticipatorias obliga a
    precisar que Codd no las plantea como un catálogo de requisitos, sino como
    criterios: varias de ellas no se cumplen íntegramente en la práctica.
  \item \textbf{Sirven como criterio de evaluación y de diseño.} Leídas como
    una lista de requisitos, las reglas describen sistemas que hoy no existen.
    Leídas como un criterio, cumplen la función opuesta: señalan con precisión
    qué sacrifica un sistema de gestión de bases de datos al apartarse del
    modelo.
\end{itemize}

\begin{conceptobox}[Una lectura crítica]
Ningún sistema relacional comercial cumple las doce reglas al pie de la letra,
  y conviene señalar por qué. La tercera es la más conflictiva: SQL incorpora
  el valor nulo como un miembro más de los dominios, lo que resuelve problemas
  prácticos, porque los datos ausentes son inevitables, a costa de introducir
  una semántica de tres valores que complica la lógica de las restricciones.
  La quinta, por su parte, cede ante consideraciones de rendimiento: la
  desnormalización y la duplicación de datos calculados se practican de forma
  habitual. La segunda supone, en fin, un control de acceso que ningún sistema
  practica de manera uniforme. Que el modelo puro y los sistemas reales
  coincidan sólo parcialmente no invalida las reglas: indica, más bien, que
  funcionan como una especificación ideal desde la cual medir la distancia
  \cite{silberschatz2019database}.
\end{conceptobox}
